\subsection{Fredholm Applications}

Let E and F be locally convex spaces. In this section, we denote by \(u \equiv v\) the congruence modulo \(\mathcal{L}^{\mathrm{f}}(\mathrm{E};\mathrm{F})\) of the elements u and v of \(\mathcal{L}(\mathrm{E};\mathrm{F})\) .

If G is a locally convex space, and if \(u'\) and \(v'\) are elements of \(\mathcal{L}(\mathrm{F};\mathrm{G})\) , the relations \(u \equiv v\) and \(u' \equiv v'\) imply the relation \(u' \circ u \equiv v' \circ v\) .

We say that an element \(w\) of \(\mathcal{L}(\mathrm{F};\mathrm{E})\) is a quasi-inverse of the element \(u\) of \(\mathcal{L}(\mathrm{E};\mathrm{F})\) if we have \(w \circ u \equiv 1_{\mathrm{E}}\) and \(u \circ w \equiv 1_{\mathrm{F}}\) .

Suppose that w is a quasi-inverse of u. If \(u_{1}\) is an element of \(\mathcal{L}(\mathrm{E};\mathrm{F})\) and \(w_{1}\) is an element of \(\mathcal{L}(\mathrm{F};\mathrm{E})\) such that \(u_{1} \equiv u\) and \(w_{1} \equiv w\) , then \(w_{1}\) is a quasi-inverse of \(u_{1}\) since \(w \circ u \equiv w_{1} \circ u_{1}\) and \(u \circ w \equiv u_{1} \circ w_{1}\) .

If w and \(w_{1}\) are quasi-inverses of u, then \(w_{1} \equiv w\) since

\[
w _ {1} = 1 _ {\mathrm{E}} \circ w _ {1} \equiv w \circ u \circ w _ {1} \equiv w \circ 1 _ {\mathrm{F}} = w.
\]

DEFINITION 1. --- Let E and F be locally convex spaces. We say that an element u of \(\mathcal{L}(\mathrm{E};\mathrm{F})\) is a Fredholm map \(^{(1)}\) if it possesses a quasi-inverse. A Fredholm map from E to E is called a Fredholm endomorphism of E.

We shall denote by \(\mathcal{F}(\mathrm{E};\mathrm{F})\) the set of Fredholm maps from E to F, and \(\mathcal{F}(\mathrm{E})\) the set of Fredholm endomorphisms of E.

Remarks. --- Let E, F, and G be locally convex spaces, \(u: E \to F\) and \(v: F \to G\) continuous linear maps.

\begin{enumerate}
\def\labelenumi{\arabic{enumi})}
\item
  Suppose that u is a Fredholm map and denote by \(u_{1}\) a quasi-inverse of u. Since u is a quasi-inverse of \(u_{1}\) , the map \(u_{1}\) is a Fredholm map.
\item
  Suppose that u and v are Fredholm maps, and let \(u_{1}\) (resp. \(v_{1}\) ) be a quasi-inverse of u (resp. v). Then \(v \circ u\) is a Fredholm map from E to G and \(u_{1} \circ v_{1}\) is a quasi-inverse of \(v \circ u\) . Indeed, one computes
\end{enumerate}

\[
\left(u _ {1} \circ v _ {1}\right) \circ (v \circ u) = u _ {1} \circ \left(v _ {1} \circ v\right) \circ u \equiv u _ {1} \circ 1 _ {\mathrm{F}} \circ u = u _ {1} \circ u \equiv 1 _ {\mathrm{E}},
\]

\[
(v \circ u) \circ (u _ {1} \circ v _ {1}) = v \circ (u \circ u _ {1}) \circ v _ {1} \equiv v \circ 1 _ {\mathrm{F}} \circ v _ {1} = v \circ v _ {1} \equiv 1 _ {\mathrm{G}}.
\]

\begin{enumerate}
\def\labelenumi{\arabic{enumi})}
\setcounter{enumi}{2}
\tightlist
\item
  Suppose that u and \(v \circ u\) are Fredholm maps, and let \(w_{1}\) be a quasi-inverse of \(v \circ u\) . Then v is a Fredholm map and \(u \circ w_{1}\) is a quasi-inverse of v.
\end{enumerate}

Indeed, \(w_{1}\) is a Fredholm operator by the first remark. Let \(u_{1}\) be a quasi-inverse of u; by the second remark, \(u \circ w_{1}\) is a Fredholm map and \((v \circ u) \circ u_{1}\) is a quasi-inverse. We have \(u \circ u_{1} \equiv 1_{F}\) , whence \(v \circ u \circ u_{1} \equiv v\) ; this proves the assertion.

\begin{enumerate}
\def\labelenumi{\arabic{enumi})}
\setcounter{enumi}{3}
\tightlist
\item
  Suppose that v and \(v \circ u\) are Fredholm maps, and let \(w_{1}\) be a quasi-inverse of \(v \circ u\) . Then u is a Fredholm operator and \(w_{1} \circ v\) is a quasi-inverse of u.
\end{enumerate}

The proof is analogous to that of the preceding remark.

Lemma 1. --- Let E and F be locally convex spaces and u a Fredholm map from E into F. The kernel and cokernel of u are finite-dimensional.

Let \(v \colon \mathrm{F} \to \mathrm{E}\) be a quasi-inverse of \(u\) . The kernel of \(u\) is contained in the image of the finite-rank linear map \(1_{\mathrm{E}} - v \circ u\) , hence is finite-dimensional. The image of \(u\) contains the kernel of the finite-rank linear map \(1_{\mathrm{F}} - u \circ v\) , hence is of finite codimension in F.

PROPOSITION 2. — Let E and F be separated locally convex spaces and let u be an element of \(\mathcal{L}(\mathrm{E};\mathrm{F})\) . The following properties are equivalent:

\begin{enumerate}
\def\labelenumi{(\roman{enumi})}
\item
  The map u is a Fredholm operator;
\item
  The map \(u\) is a strict morphism, its kernel is finite-dimensional, its image is closed and of finite codimension in F;
\item
  There exist closed vector subspaces \(E_{1}\) of E and \(F_{1}\) of F, both of finite codimension, such that u induces, by passing to subspaces, an isomorphism from \(E_{1}\) onto \(F_{1}\) ;
\item
   There exist decompositions into topological direct sums \(E = E_{1} \oplus E_{2}\) and \(F = F_{1} \oplus F_{2}\) such that \(E_{2}\) and \(F_{2}\) are finite-dimensional, that u vanishes on \(E_{2}\) and defines by passing to subspaces an isomorphism of \(E_{1}\) onto \(F_{1}\) .
\item
  \(\Longrightarrow\) (iii): Let v be a quasi-inverse of u, \(E_{1}\) the kernel of \(1_{E}-v\circ u\) and \(F_{1}\) that of \(1_{F}-u\circ v\) . Since the linear maps \(1_{E}-v\circ u\) and \(1_{F}-u\circ v\) are continuous and of finite rank, \(E_{1}\) and \(F_{1}\) are closed vector subspaces of finite codimension of E and F, respectively. Let \(x\in E_{1}\) . We have
\end{enumerate}

\[
(1 _ {\mathrm{F}} - u \circ v) (u (x)) = u ((1 _ {\mathrm{E}} - v \circ u) (x)) = u (0) = 0,
\]

whence \(u(x) \in \mathrm{F}_1\) . We therefore have \(u(\mathrm{E}_1) \subset \mathrm{F}_1\) ; we likewise have \(v(\mathrm{F}_1) \subset \mathrm{E}_1\) . The continuous linear maps \(u_1: \mathrm{E}_1 \to \mathrm{F}_1\) and \(v_1: \mathrm{F}_1 \to \mathrm{E}_1\) deduced from \(u\) and \(v\) are then mutually inverse isomorphisms, since \(v \circ u\) and \(1_{\mathrm{E}}\) (resp. \(u \circ v\) and \(1_{\mathrm{F}}\) ) coincide on \(\mathrm{E}_1\) (resp. on \(\mathrm{F}_1\) ).

\begin{enumerate}
\def\labelenumi{(\roman{enumi})}
\setcounter{enumi}{2}
\item
  \(\Longrightarrow\) (ii): Let \(E_{1}\) and \(F_{1}\) satisfy the hypothesis of (iii). We have \(\mathrm{E}_{1}\cap\mathrm{Ker}(u)=\{0\}\) and \(\mathrm{F}_{1}\subset\mathrm{Im}(u)\) , hence \(\mathrm{Ker}(u)\) is finite-dimensional and \(\mathrm{Im}(u)\) is closed and of finite codimension in F. It follows from prop. 1 of III, p. 39 that the map u is a strict morphism.
\item
   \(\Longrightarrow\) (iv): Suppose condition (ii) is satisfied. The closed vector subspace \(E_{2} = \operatorname{Ker}(u)\) of E is finite-dimensional, and there exists a vector subspace \(E_{1}\) of E topologically complementary to \(E_{2}\) (EVT, II, p. 27, cor. 2). The vector subspace \(F_{1} = \operatorname{Im}(u)\) is closed and of finite codimension, and it admits a topological complement \(F_{2}\) in F. By prop. 1 of III, p. 39, the map \(u|E_{1}\) is a strict morphism, hence u induces an isomorphism of \(E_{1}\) onto \(F_{1}\) .
\item
   \(\Longrightarrow\) (i): Under the hypotheses of (iv), the linear map from F to E that coincides with \(u_1^{-1}\) on \(\mathrm{F}_1\) and is zero on \(\mathrm{F}_2\) is a quasi-inverse of \(u\) .
\end{enumerate}

Remark 5. --- Let E, F be separated locally convex spaces and \(u\colon\mathrm{E}\to\mathrm{F}\) a Fredholm operator. If \(u\) is bijective, then \(u\) is an isomorphism (indeed, \(u\) is a strict morphism by Prop. 2, (ii)).

